\makeatletter
\@ifundefined{abstract}{%
  \newenvironment{abstract}{}{}
}{}
\makeatother

\begin{abstract}
The challenge of selecting a container network plugin in Kubernetes is critical for modern cloud-native architectures. In this paper, we study the performance and security trade-offs of CNI implementations under Zero Trust network policies. We design a comprehensive evaluation framework and a Multi-Criteria Decision Analysis model. Our results show that eBPF-based CNI plugins reduce CPU overhead by 15\% over iptables-based implementations under security rules. These findings demonstrate that selecting the appropriate CNI plugin mitigates resource overprovisioning.
\end{abstract}

\chapter{Estado del Arte} \label{sec:EstadoArte}
% \section{Related Work}

El diseño y la optimización del plano de datos en la orquestación de contenedores constituye un área de activa investigación. En este capítulo se analiza y contrasta la literatura científica y técnica reciente referente a la evaluación de complementos de red en Kubernetes, la aplicación de políticas de seguridad bajo esquemas de confianza cero y el uso de modelos de decisión para la selección de infraestructura.

\section{Evaluación Comparativa de Rendimiento de CNIs en la Literatura}

La evaluación experimental de complementos CNI representa un campo recurrente debido al impacto directo de la red en la eficiencia de las aplicaciones. Diversos autores han examinado el comportamiento de complementos populares como Calico, Cilium, Flannel y Antrea bajo diversos escenarios.

Un estudio pionero desarrollado por Kang et al.~\cite{ref7} evaluó el rendimiento de múltiples plugins CNI en entornos de edge computing con recursos limitados. Los autores analizaron variables como throughput del Protocolo de Control de Transmisión (TCP), la latencia de red y el uso de CPU, concluyendo que Flannel destaca por su baja sobrecarga en escenarios con limitaciones de cómputo. Sin embargo (\emph{however}), esta investigación no contempló el impacto del cifrado de red o de las políticas de seguridad en la degradación de las métricas. 

En contraste (\emph{in contrast}), Koukis et al.~\cite{ref8} extendieron la comparación de rendimiento a entornos multi-nodo utilizando herramientas de benchmarking automatizadas. Sus resultados revelaron que los plugins basados en enrutamiento nativo superan sistemáticamente a las soluciones superpuestas (overlay) en throughput. No obstante, reconocen que el aislamiento lógico que proveen los overlays simplifica la operación en nubes públicas, estableciendo un compromiso técnico (\emph{trade-off}) entre rendimiento y simplicidad operativa.

Por su parte, el trabajo de Pfaff et al.~\cite{ref22} sobre Open vSwitch (OVS) demostró que las arquitecturas de conmutación basadas en software de alto rendimiento disminuyen la sobrecarga de procesamiento en comparación con (\emph{compared to}) soluciones basadas en el filtrado tradicional del kernel. Esta fundamentación teórica y práctica justifica la arquitectura adoptada por Antrea, que utiliza OVS para optimizar la conmutación interna del clúster.

Adicionalmente, investigaciones recientes enfocadas en el kernel de Linux mediante eBPF (extended Berkeley Packet Filter), como las presentadas por Zhang et al.~\cite{ref27}, evidencian que compilar filtros y políticas directamente en instrucciones del Protocolo de Internet (IP) del kernel mitiga drásticamente el uso de CPU. A diferencia de (\emph{unlike}) iptables, cuya complejidad de búsqueda es lineal respecto al número de reglas, el uso de mapas hash en eBPF mantiene un tiempo de procesamiento constante. Esto confiere a Cilium y Calico (en su modo eBPF) una ventaja competitiva significativa en clústeres de gran escala.

En resumen (\emph{in summary}), la literatura científica reporta una clara divergencia en los resultados de rendimiento de red según el dataplane empleado, donde el uso de eBPF representa la tendencia predominante para clústeres de gran volumen, mientras que Flannel sigue siendo relevante en escenarios de escasos recursos debido a su simplicidad.

\section{Seguridad y Políticas de Red en la Literatura Académica}

El control de acceso de red en entornos nativos de la nube ha evolucionado desde perímetros estáticos hacia modelos dinámicos orientados a la identidad del servicio. La literatura examina de forma crítica las implicaciones de implementar políticas de red estrictas.

La arquitectura de confianza cero (Zero Trust Architecture) formalizada por el Instituto Nacional de Estándares y Tecnología (NIST) en su Publicación Especial (SP) 800-207 por Rose et al.~\cite{ref24} establece que no se debe conceder confianza implícita a ninguna carga de trabajo por su ubicación física o lógica. Este paradigma exige la verificación constante y la aplicación de políticas de mínimo privilegio a nivel de microservicio, un lineamiento que en Kubernetes se delega en las NetworkPolicies~\cite{ref25}.

Sin embargo (\emph{however}), la implementación de estas políticas distribuidas introduce un overhead que los autores clásicos suelen omitir. Chandramouli et al.~\cite{ref26} evaluaron la seguridad en microservicios utilizando arquitecturas de malla de servicios (Service Mesh). Aunque (\emph{although}) confirmaron que el uso de mTLS y proxies L7 garantiza un control exhaustivo, reportaron un incremento notable en la latencia y en el consumo de CPU. Esto plantea una limitación común (\emph{common limitation}) para organizaciones que buscan proteger su tráfico sin sobredimensionar sus nodos.

En el ámbito del monitoreo de red, Gomes et al.~\cite{ref6} realizaron un mapeo sistemático de la observabilidad en microservicios. Los autores identificaron que la exportación de telemetría de flujos correlacionada con namespaces y etiquetas lógicas es fundamental para la auditoría de seguridad. No obstante, alertaron que el volumen de datos recolectados mediante protocolos como la Exportación de Información de Flujo de IP (IPFIX) genera una sobrecarga de almacenamiento y procesamiento que debe ser evaluada empíricamente~\cite{ref23}.

De esta manera, las investigaciones previas coinciden en que la seguridad Zero Trust es indispensable, pero su coste en rendimiento y consumo de recursos sigue siendo una variable crítica insuficientemente cuantificada en la literatura actual.

\section{Modelos de Decisión Multicriterio en Infraestructura Cloud}

La selección de infraestructura en la nube y complementos tecnológicos representa un desafío complejo debido a la presencia de múltiples criterios contrapuestos (e.g., latencia baja vs. consumo de CPU bajo). La literatura científica activa aborda este problema mediante modelos formales.

Los modelos de Análisis de Decisión Multicriterio (MCDA) han sido ampliamente utilizados en la ingeniería de sistemas para la selección de alternativas. Investigaciones previas han empleado metodologías como el Proceso de Jerarquía Analítica (AHP) y la Técnica de Orden de Preferencia por Similitud con la Solución Ideal (TOPSIS) para clasificar servicios en la nube en función del costo, disponibilidad y rendimiento.

A pesar de (\emph{despite}) la madurez de estos modelos de decisión, su aplicación específica al ecosistema de red de Kubernetes es escasa. La mayoría de las decisiones de selección de CNI en la industria se basan en recomendaciones empíricas o heurísticas simples, ignorando el compromiso estructurado (\emph{trade-off}) que existe entre rendimiento, seguridad y huella de recursos del clúster.

\subsection{Brecha de Investigación (Research Gap)}

Al integrar las tres vertientes de la literatura de red examinadas, se evidencia una clara brecha de investigación (\emph{research gap}). A pesar de los valiosos estudios comparativos de rendimiento de CNI~\cite{ref7,ref8} y del análisis conceptual de la seguridad Zero Trust~\cite{ref24,ref26}, la literatura actual carece de análisis experimentales reales que midan de forma simultánea el rendimiento de red (QoS), el consumo energético del clúster y el coste de CPU derivado de la validación de políticas de seguridad automatizadas en nubes públicas.

Esta carencia de datos integrales representa una limitación metodológica (\emph{methodological limitation}) crítica que obliga a los equipos de ingeniería a tomar decisiones de infraestructura subóptimas o a sobredimensionar preventivamente los clústeres. Por consiguiente, esta tesis cubre dicho vacío al proveer un modelo de evaluación cuantitativo y un framework de decisión multicriterio estructurado para la selección informada de complementos CNI bajo políticas Zero Trust.